% Figure: net lift of a fixed-volume gas balloon against altitude. Both the
% surrounding air density and (for a sealed envelope at ambient pressure) the lift
% gas density fall with altitude in the same ratio, so the net lift of a constant-mass
% airship declines with height and crosses zero at the pressure (or ceiling) altitude
% where the displaced-air weight equals the craft weight. No em-dashes.
\begin{figure}[htbp]
\centering
\begin{tikzpicture}
\begin{axis}[
  width=10cm, height=6.6cm,
  xlabel={altitude $z$ (km)},
  ylabel={net lift $L_{\text{net}}$ (kN)},
  xmin=0, xmax=8,
  ymin=-4, ymax=18,
  axis lines=left,
  grid=major,
  grid style={enggray!25},
  tick label style={font=\scriptsize},
  label style={font=\small},
  legend style={font=\scriptsize, at={(0.97,0.97)}, anchor=north east, draw=enggray!50},
]
% Net lift L_net(z) = g*V*(rho_air(z) - rho_gas(z)) - W_struct
% Model: rho_air(z) = rho0 * exp(-z/H), helium scales the same way for a vented
% (pressure-equalised) envelope, so buoyant force per unit displaced volume falls
% exponentially. Use rho0=1.225, rho_he0=0.169, V=12000 m^3, H=8.4 km,
% W_struct = 90 kN (constant). Net lift in kN:
%   L(z) = 9.81*12000*(1.225-0.169)*exp(-z/8.4)/1000 - 9.0   [kN], tuned to cross ~0
% gross buoyant-minus-gas lift at sea level ~ 124 kN scaled down for a legible plot
\addplot[engblue, very thick, samples=80, domain=0:8]
  ({x}, {124.3*exp(-x/8.4) - 110});
\addlegendentry{net lift $L_{\text{net}}(z)$}
% zero line
\addplot[engred, dashed, thick, samples=2, domain=0:8] ({x},{0});
\addlegendentry{neutral buoyancy}
% mark the pressure altitude where L_net=0
% 124.3*exp(-z/8.4)=110  ->  z = -8.4*ln(110/124.3) ~ 1.02 km
\addplot[only marks, mark=*, engblack, mark size=1.6pt] coordinates {(1.02,0)};
\node[engblack, font=\scriptsize, anchor=south west] at (axis cs:1.02,0.3) {pressure ceiling};
\end{axis}
\end{tikzpicture}
\caption[Net lift of an airship against altitude]{Net lift of a
constant-mass, pressure-equalised gas airship against altitude. The buoyant
force is the weight of displaced air, and the displaced-air weight falls
roughly exponentially as the surrounding density drops with height (the
barometric law); the lift-gas weight falls in the same ratio for a vented
envelope, so the gas contributes a fixed fraction and the net lift tracks the
air density. Net lift crosses zero at the pressure ceiling, the altitude at
which the displaced-air weight has fallen to the craft weight. To climb
higher the craft must shed ballast or compress and store gas; to descend it
must valve gas or take on ballast. The curve is the static-equilibrium
backbone of every lighter-than-air design.}
\label{fig:vol03ch10:airship-equilibrium}
\end{figure}
